\documentclass[11pt,a4paper]{ctexart}
\usepackage[student]{ipara}

\pagestyle{fancy}
\fancyhf{}
\lhead{\small 专题讲义：函数的对称性与周期性}
\rhead{\small 姓名：\blank{2.5cm}}
\cfoot{\small 第 \thepage\ 页\quad 共 \zpageref{LastPage} 页}
\renewcommand{\headrulewidth}{0.4pt}
\renewcommand{\footrulewidth}{0pt}

\begin{document}

\begin{center}
  {\LARGE\bfseries 专题讲义：函数的对称性与周期性}\par
  \vspace{0.6em}
  {\small 姓名：\blank{2.5cm} \quad 日期：\blank{2.5cm} \quad 课次：第 X 讲}
\end{center}
\vspace{0.6em}

\begin{tcolorbox}[
  enhanced,
  colback=white,
  colframe=black,
  boxrule=0.8pt,
  arc=0pt,
  left=10pt, right=10pt, top=8pt, bottom=8pt,
  title={\small\bfseries 【核心教学哲学】}
]
  {\small
  \textbf{先在图上看见平移，再用代数把平移证明出来。}\par
  对称是局部的翻折，周期是整体的平移；\textbf{两次合适的翻折，往往就等于一次平移}。\par
  周期函数不是“麻烦的函数”，而是“信息高度重复的函数”。一旦发现周期性，就等于把整条实数轴上的问题，压缩到了一个周期内。
  }
\end{tcolorbox}

\section{核心直观：为什么“两次对称”会复合出平移？}

很多同学卡在周期性公式，本质上是因为没有看见几何过程。建议采用\textbf{“上图下式”}的板书与思考方式：

\subsection{1. 两条平行对称轴：轴对称 + 轴对称 = 平移}

画一小段不规则曲线，先后关于直线 $x=a$ 与 $x=b$ 翻折：

\begin{center}
\begin{tikzpicture}[scale=0.85, >=stealth]
  \draw[->, gray!60, thin] (-0.5,0) -- (7.5,0) node[right, black] {$x$};
  \draw[->, gray!60, thin] (0,-0.8) -- (0,2.5) node[above, black] {$y$};
  \draw[dashed, thick] (2,-0.6) -- (2,2.3) node[above] {$x=a$};
  \draw[dashed, thick] (4,-0.6) -- (4,2.3) node[above] {$x=b$};
  
  % C0: x in [0,2]
  \draw[very thick] plot[smooth, tension=0.7] coordinates {(0, 0.5) (1, 1.9) (2, 0.8)};
  \node[above] at (1, 1.9) {$C_0$};
  \fill (1, 1.9) circle (1.5pt) node[above left] {$P(x,y)$};
  
  % C1: x in [2,4] (关于 x=a=2 轴对称)
  \draw[thick, densely dotted] plot[smooth, tension=0.7] coordinates {(2, 0.8) (3, 1.9) (4, 0.5)};
  \node[above] at (3, 1.9) {$C_1$};
  \fill (3, 1.9) circle (1.5pt) node[above right] {$P_1(2a-x, y)$};
  \draw[->, bend left=25, darkgray] (1.0, 2.1) to node[above, font=\small] {$x=a$ 翻折} (3.0, 2.1);
  
  % C2: x in [4,6] (关于 x=b=4 轴对称)
  \draw[very thick] plot[smooth, tension=0.7] coordinates {(4, 0.5) (5, 1.9) (6, 0.8)};
  \node[above] at (5, 1.9) {$C_2$};
  \fill (5, 1.9) circle (1.5pt) node[above right] {$P_2(x+2(b-a), y)$};
  \draw[->, bend left=25, darkgray] (3.0, 2.1) to node[above, font=\small] {$x=b$ 翻折} (5.0, 2.1);
  
  \draw[->, double, thick] (1.0, -0.4) -- (5.0, -0.4) node[midway, below, font=\small\bfseries] {无缝连贯：整体向右平移 $2(b-a)$};
\end{tikzpicture}
\end{center}

\[
x \overset{x=a}{\longmapsto} 2a-x \overset{x=b}{\longmapsto} 2b-(2a-x) = x+2(b-a)
\]

\textbf{结论}：翻折两次 $=$ \blank{4cm}。因此 \blank{4cm} 是 $f(x)$ 的一个周期。

\subsection{2. 两个等高对称中心：中心对称 + 中心对称 = 平移}

设两个对称中心为 $A(a,m)$ 与 $B(b,m)$：

\begin{center}
\begin{tikzpicture}[scale=0.85, >=stealth]
  \draw[->, gray!60, thin] (-0.5,0) -- (7.5,0) node[right, black] {$x$};
  \draw[->, gray!60, thin] (0,-0.8) -- (0,2.5) node[above, black] {$y$};
  \draw[gray!40, dashed] (-0.5, 0.8) -- (7.5, 0.8) node[right, black, font=\small] {$y=m$};
  
  \fill (2.0, 0.8) circle (2pt) node[below right] {$A(a,m)$};
  \fill (4.0, 0.8) circle (2pt) node[below right] {$B(b,m)$};
  
  % C0: x in [0,2]
  \draw[very thick] plot[smooth, tension=0.7] coordinates {(0, 0.8) (1, 1.9) (2, 0.8)};
  \node[above] at (1, 1.9) {$C_0$};
  \fill (1, 1.9) circle (1.5pt) node[above left] {$P(x,y)$};
  
  % C1: x in [2,4] (绕 A(2,0.8) 旋转 180 度)
  \draw[thick, densely dotted] plot[smooth, tension=0.7] coordinates {(2, 0.8) (3, -0.3) (4, 0.8)};
  \node[below] at (3, -0.3) {$C_1$};
  \fill (3, -0.3) circle (1.5pt) node[below right] {$P_1(2a-x, 2m-y)$};
  \draw[dashed, gray] (1, 1.9) -- (3, -0.3);
  
  % C2: x in [4,6] (绕 B(4,0.8) 旋转 180 度)
  \draw[very thick] plot[smooth, tension=0.7] coordinates {(4, 0.8) (5, 1.9) (6, 0.8)};
  \node[above] at (5, 1.9) {$C_2$};
  \fill (5, 1.9) circle (1.5pt) node[above right] {$P_2(x+2(b-a), y)$};
  \draw[dashed, gray] (3, -0.3) -- (5, 1.9);
  
  \draw[->, double, thick] (1.0, 2.2) -- (5.0, 2.2) node[midway, above, font=\small\bfseries] {无缝连贯：中心对称两次 $=$ 向右平移 $2(b-a)$};
\end{tikzpicture}
\end{center}

\textbf{结论}：绕点 $A$ 旋转 $180^\circ$，再绕点 $B$ 旋转 $180^\circ$，最终效果为向右平移 \blank{3cm}。

\subsection{3. 一条对称轴与一个对称中心：轴对称 + 中心对称 = 反相平移（两步恢复）}

设对称轴为 $x=a$，对称中心为 $B(b,m)$：

\begin{center}
\begin{tikzpicture}[scale=0.8, >=stealth]
  \draw[->, gray!60, thin] (-0.5,0) -- (10.5,0) node[right, black] {$x$};
  \draw[->, gray!60, thin] (0,-1.5) -- (0,2.5) node[above, black] {$y$};
  \draw[dashed, thick] (2,-1.2) -- (2,2.3) node[above] {$x=a$};
  \draw[gray!40, dashed] (-0.5, 0.8) -- (10.5, 0.8) node[right, black, font=\small] {$y=m$};
  \fill (4.0, 0.8) circle (2pt) node[below right] {$B(b,m)$};
  
  % C0: x in [0,2]
  \draw[very thick] plot[smooth, tension=0.7] coordinates {(0, 0.8) (1, 2.0) (2, 0.8)};
  \node[above] at (1, 2.0) {$C_0$ (正位)};
  
  % C1: x in [2,4]
  \draw[thick, densely dotted] plot[smooth, tension=0.7] coordinates {(2, 0.8) (3, 2.0) (4, 0.8)};
  \node[above] at (3, 2.0) {$C_1$};
  
  % C2: x in [4,6]
  \draw[thick, densely dotted] plot[smooth, tension=0.7] coordinates {(4, 0.8) (5, -0.4) (6, 0.8)};
  \node[below] at (5, -0.4) {$C_2$ (颠倒: $f(x+T)=2m-f(x)$)};
  
  % C3: x in [6,8]
  \draw[thick, densely dotted] plot[smooth, tension=0.7] coordinates {(6, 0.8) (7, -0.4) (8, 0.8)};
  \node[below] at (7, -0.4) {$C_3$};
  
  % C4: x in [8,10]
  \draw[very thick] plot[smooth, tension=0.7] coordinates {(8, 0.8) (9, 2.0) (10, 0.8)};
  \node[above] at (9, 2.0) {$C_4$ (恢复正立)};
  
  \draw[->, double, thick] (1.0, 2.3) -- (9.0, 2.3) node[midway, above, font=\small\bfseries] {无缝连续曲线：两步恢复周期为 $4|b-a|$};
\end{tikzpicture}
\end{center}

第一次复合后，图像虽平移，但上下颠倒 $f(x+2(b-a)) = 2m - f(x)$（反相平移）；再进行一次相同变化，图像才完全恢复，因此周期为 \blank{3cm}。

\newpage

\section{三大“对称复合模型”代数精讲}

不要凭感觉换元！代数推导的核心目标：\textbf{制造 $f(x+T)$}。

可以在草稿纸上写出箭号推导流程：
\[
x \xrightarrow{x=a \text{ 对称}} 2a-x \xrightarrow{x=b \text{ 对称}} 2b-(2a-x) = x+2(b-a)
\]

\subsection{模型一：两条对称轴 $\implies$ 周期}
若 $f(x)$ 关于 $x=a$ 和 $x=b$ 均对称，则：
\[
f(2a-x) = f(x), \qquad f(2b-x) = f(x).
\]
推导：$f\bigl(x+2(b-a)\bigr) = f\bigl(2b-(2a-x)\bigr) = f(2a-x) = f(x)$。
\[
\boxed{2|b-a| \text{ 是 } f(x) \text{ 的一个周期}}
\]

\subsection{模型二：两个中心对称 $\implies$ 周期或斜向重复}
若 $f(x)$ 关于 $A(a,m)$ 和 $B(b,n)$ 中心对称，则：
\[
f(2a-x) = 2m - f(x), \qquad f(2b-x) = 2n - f(x).
\]
推导：$f\bigl(x+2(b-a)\bigr) = 2n - f(2a-x) = 2n - (2m - f(x)) = f(x) + 2(n-m)$。
\begin{itemize}[itemsep=0.3em]
  \item \textbf{若 $m=n$（等高中心）}：$f\bigl(x+2(b-a)\bigr) = f(x) \implies \boxed{2|b-a| \text{ 是一个周期}}$。
  \item \textbf{若 $m \neq n$（不等高中心）}：$f(x+T) = f(x) + C \implies$ \textbf{准周期式关系（斜向重复，非周期函数）}。
\end{itemize}

\subsection{模型三：一条对称轴和一个对称中心 $\implies$ 四倍周期}
若 $f(x)$ 关于 $x=a$ 轴对称，关于 $B(b,m)$ 中心对称：
\[
f(2a-x) = f(x), \qquad f(2b-x) = 2m - f(x).
\]
一次复合得反相平移：$f\bigl(x+2(b-a)\bigr) = 2m - f(2a-x) = 2m - f(x)$。\par
再平移一次：$f\bigl(x+4(b-a)\bigr) = 2m - f\bigl(x+2(b-a)\bigr) = 2m - (2m - f(x)) = f(x)$。
\[
\boxed{4|b-a| \text{ 是 } f(x) \text{ 的一个周期}}
\]

\section{学生必须条件反射的四个表达式}

\begin{enumerate}[label=\arabic*. , leftmargin=2em, itemsep=0.4em]
  \item \textbf{轴对称}：$f(a+x) = f(a-x) \iff f(2a-x) = f(x) \iff$ 图像关于 \blank{2cm} 对称。
  \item \textbf{中心对称}：$f(a+x) + f(a-x) = 2b \iff f(2a-x) = 2b - f(x) \iff$ 图像关于点 \blank{2cm} 中心对称。特别地，$f(a+x) = -f(a-x) \iff$ 中心为 $(a,0)$。
  \item \textbf{直接周期}：$f(x+T) = f(x) \iff T \neq 0$ 是一个周期。
  \item \textbf{反周期/反相平移}：$f(x+T) = -f(x)$ 或 $f(x+T) = 2m - f(x) \iff$ \blank{2cm} 是一个周期。
\end{enumerate}

\section{为什么看到周期函数应该高兴？（降维打击五大优势）}

\begin{enumerate}[label=\arabic*. , leftmargin=2em, itemsep=0.4em]
  \item \textbf{无限压缩到有限}：只需研究长度为 $T$ 的区间（如 $[0,T)$）。
  \item \textbf{大数化小}：若 $T=6$，则 $f(2026) = f(337 \times 6 + 4) = f(4)$。
  \item \textbf{零点成组复制}：若 $f(x_0)=0$，则 $f(x_0+nT)=0 \quad (n \in \mathbb{Z})$。
  \item \textbf{最值与单调性重复}：只需分析一个周期内的最值点与单调区间。
  \item \textbf{定积分搬移}：$\int_{a}^{a+T} f(x) \, dx = \int_{0}^{T} f(x) \, dx$。若 $f(x+T) = -f(x)$，则 $\int_{a}^{a+2T} f(x) \, dx = 0$。
\end{enumerate}

\newpage

\section{固定四步法与典例演练}

\mindsetcard{固定四步法}{解题规范路线图}{
  \textbf{第一步}：几何翻译——将轴/中心对称翻译成函数关系式；\\
  \textbf{第二步}：轨迹计算——先计算连续两次变换后的横坐标变化 $x \to 2a-x \to x+2(b-a)$；\\
  \textbf{第三步}：连续代换——连续代入对称式制造 $f(x+T)$ 或 $f(x+T)=C-f(x)$；\\
  \textbf{第四步}：下结论——准确写出周期 $2|b-a|$ 或 $4|b-a|$（不随意写“最小”）。
}

\studentproblem{1}
已知定义在 $\mathbb{R}$ 上的函数 $f(x)$ 满足 $f(1+x) = f(1-x)$，且 $f(3+x) = -f(3-x)$。
\begin{enumerate}[label=(\arabic*)]
  \item 求证：$f(x)$ 是周期函数，并求出其一个周期；
  \item 若 $f(1) = 2$，求 $f(2025) + f(2026)$ 的值。
\end{enumerate}

\answerblank[8.5cm]{解：}

\studentthink{
  \item 我是否准确辨析了 $f(1+x)=f(1-x)$（对称轴）与 $f(3+x)=-f(3-x)$（对称中心）？
  \item 我求出的周期是 $4$ 还是 $8$？是否经过了“两步恢复”验证？
}

\studentproblem{2}
设 $f(x)$ 是定义在 $\mathbb{R}$ 上的连续函数，满足 $f(x+2) = -f(x)$。
\begin{enumerate}[label=(\arabic*)]
  \item 求 $f(x)$ 的一个周期；
  \item 若 $f(x)$ 在 $[0,2]$ 上的零点个数为 $3$，求 $f(x)$ 在 $[0,10]$ 上的零点个数。
\end{enumerate}

\answerblank[8.5cm]{解：}

\studentthink{
  \item $f(x+2)=-f(x)$ 属于反周期关系，周期应当是 $2$ 还是 $4$？
  \item 区间 $[0,10]$ 包含多少个周期？端点零点是否重复计算？
}

\newpage

\section{课堂错因大起底与自查}

\vspace{0.5em}
\noindent\textbf{本讲我的主要问题集中在（请打勾）：}
\begin{itemize}[itemsep=0.4em]
  \item \checkbox 混淆对称与周期：把关于 $x=a$ 轴对称误写成 $f(x+a)=f(x)$。
  \item \checkbox 忽视频率与周期倍数：一轴一中心误算成 $2|b-a|$，漏掉上下颠倒后的二次恢复。
  \item \checkbox 严谨性缺失：未证明前直接断言为“最小正周期”。
  \item \checkbox 大数化小取模计算失误：对负数或大整数取模算错余数。
\end{itemize}

\vspace{1em}
\noindent\textbf{本讲最需要记牢的一句口诀：}
\vspace{0.5em}

\begin{tcolorbox}[
  enhanced,
  colback=gray!5,
  colframe=black,
  boxrule=0.8pt,
  arc=0pt,
  center title
]
  \centering\bfseries\large
  一轴一翻，两轴一移；两心一移，轴心反相；反相两次，必出周期。
\end{tcolorbox}

\vspace{1em}
\noindent\textbf{我的课后一句话总结：}
\vspace{0.5em}

\blank{14cm}

\end{document}
